%
% presence-heartbeat.tex
%
% Z Specification for the heartbeat write of a presence row (PL-1704, DES-097,
% DES-100) and for the plan write that shares the row.
%
% Companion to nats-relay.tex: that specification models the sessions bucket
% as SessionStore (sessions : SESSIONKEY \pinj KVVALUE) and says nothing about
% two writers racing on one row.  This one models the race, the
% compare-and-set that closes it, the activity stamp the heartbeat carries, and
% the guard on an automatic plan.  It is kept apart because every component
% added to the composed state of nats-relay.tex multiplies a state space that
% already holds 831,363 states.
%
% Type-check: fuzz -t presence-heartbeat.tex
% Model-check: see the end of this file.
%

\documentclass[a4paper,10pt,fleqn]{article}
\usepackage[margin=1in]{geometry}
\usepackage{fuzz}
\usepackage[colorlinks=true,linkcolor=blue,citecolor=blue,urlcolor=blue]{hyperref}

\begin{document}

\title{Biff Presence Heartbeat: A Z Specification}
\author{Punt Labs}
\date{November 2026}
\maketitle

%%----------------------------------------------------------------------
\section{Introduction}
%%----------------------------------------------------------------------

A session's presence row lives in one key of the sessions bucket.  Two
writers use it.  The heartbeat (\texttt{NatsRelay.heartbeat}) reads the row,
moves \texttt{last\_active} forward, and writes the row back.  Everything
else (\texttt{plan}, \texttt{mesg}, a tty rename) reads the row, changes one
field, and writes it back.

The heartbeat's read and its write are two round trips.  A plan saved
between them is in the bucket when the heartbeat writes, and the heartbeat
writes the row it read: the plan is gone.  The loss is silent, and it is
the heartbeat's own tick that causes it.  PL-1704 is a report of exactly this.

The bucket gives every write a revision, and \texttt{kv.update} writes only
if the revision is the one the caller names.  The heartbeat now names the
revision it read.  A write that finds the row changed fails, and the
heartbeat reads again; after three failures it gives up until the next tick.
One skipped tick is harmless: the row lives for three days.

Two things ride on the same write.  The plugin hook touches a per-session
\texttt{.active} file after every tool call, and the heartbeat passes that
file's modification time as \texttt{at} (DES-097, DES-100): if a stamp exists,
\texttt{last\_tool\_at} becomes the later of its old value and \texttt{at}; if
none exists it is left alone; it never moves back.  Without it a session that
never calls a biff tool reads idle for as long as its row is stale (a row ten
minutes stale and a stamp one minute old read idle for 600 seconds).  And a
plan written with \texttt{source = auto} never replaces a non-empty plan
written with \texttt{source = manual}.

This specification states the properties and checks each twice: against the
shipped behaviour, exhaustively, and against the behaviour it replaced, where
ProB finds the counterexample.

%%----------------------------------------------------------------------
\section{Basic Types}
%%----------------------------------------------------------------------

\begin{zed}
  ZBOOL ::= ztrue | zfalse
\end{zed}

\begin{zed}
  FIELD ::= fplan | fmesg
\end{zed}

\begin{zed}
  SRC ::= srcManual | srcAuto
\end{zed}

$FIELD$ stands for every field of the row that a writer other than the
heartbeat owns.  Two of them are enough to say that the heartbeat preserves
\emph{every} such field and not one in particular.  $SRC$ is the
\texttt{plan\_source} of a plan.

\begin{axdef}
  maxRev : \nat \\
  maxBeat : \nat \\
  maxTries : \nat \\
  maxNow : \nat
  \where
  maxRev = 3 \\
  maxBeat = 1 \\
  maxTries = 2 \\
  maxNow = 2
\end{axdef}

$maxRev$, $maxBeat$ and $maxNow$ bound the revision, the heartbeat counter and
the clock so that ProB can enumerate them; $maxTries$ is the retry limit of
\texttt{heartbeat}, here two where the code has three (the argument does not
depend on the number).  Field values are the numbers $0$ (empty) and $1$.

\begin{zed}
  HbPhase ::= hbIdle | hbRead
\end{zed}

%%----------------------------------------------------------------------
\section{State}
%%----------------------------------------------------------------------

\begin{schema}{Row}
  fields : FIELD \fun (0 \upto 1) \\
  planSource : SRC \\
  lastTool : 0 \upto maxNow \\
  beat : 0 \upto maxBeat \\
  curRev : 0 \upto maxRev
\end{schema}

The row in the bucket: the fields the other writers own, the source of the
plan, the idle clock ($last\_tool\_at$), the heartbeat counter
($last\_active$), and the revision the bucket assigned to the last write.

\begin{schema}{Env}
  now : 0 \upto maxNow \\
  stamp : 0 \upto maxNow \\
  hasStamp : ZBOOL
\end{schema}

The world outside the bucket: the clock, and the \texttt{.active} file's
modification time, which exists once a tool call has touched it.

\begin{schema}{Heartbeat}
  phase : HbPhase \\
  seenFields : FIELD \fun (0 \upto 1) \\
  seenSource : SRC \\
  seenTool : 0 \upto maxNow \\
  seenRev : 0 \upto maxRev \\
  seenAt : 0 \upto maxNow \\
  seenHas : ZBOOL \\
  tries : 0 \upto maxTries
\end{schema}

The heartbeat in flight.  After its read, $seenFields$, $seenSource$,
$seenTool$ and $seenRev$ are the row and the revision it holds, and $seenAt$
and $seenHas$ the stamp it read at the start of its tick; $tries$ counts the
failed writes of this tick.

\begin{schema}{Intent}
  intended : FIELD \fun (0 \upto 1) \\
  intendedSource : SRC \\
  peak : 0 \upto maxNow \\
  manualLive : ZBOOL \\
  manualPlan : 0 \upto 1 \\
  beatHas : ZBOOL \\
  beatAt : 0 \upto maxNow
\end{schema}

Ghosts.  $intended$ and $intendedSource$ are what the other writers last wrote,
in the order they wrote it.  $peak$ is the greatest idle-clock value the row
has held.  $manualLive$ says a non-empty manual plan was written and not since
replaced by a manual write, and $manualPlan$ is its value.  $beatHas$ and
$beatAt$ are the stamp the last committed heartbeat carried.  None of these is
in the bucket and the code has no such variable; they exist so that a lost
write, a clock that went back and a replaced manual plan are states and not
arguments.

\begin{schema}{PresenceState}
  Row \\
  Env \\
  Heartbeat \\
  Intent
\end{schema}

\begin{schema}{PresenceInv}
  PresenceState
  \where
  fields = intended \\
  planSource = intendedSource \\
  lastTool = peak \\
  beatHas = ztrue \implies lastTool \geq beatAt \\
  manualLive = ztrue \implies
    ( planSource = srcManual \land fields~fplan = manualPlan )
\end{schema}

The properties, in order.  (1) The row holds, for every field and for the
plan's source, the last value a writer other than the heartbeat gave it: the
heartbeat never overwrites a plan or any other field written after its read.
(2) The idle clock never moves backwards: the row's $lastTool$ is its peak.
(3) A committed heartbeat leaves the idle clock at least as fresh as the stamp
it carried: idle is at most the age of the stamp.  (4) An automatic plan never
replaces a non-empty manual one.

\begin{schema}{Invariant}
  PresenceInv
\end{schema}

\begin{schema}{Init}
  PresenceState'
  \where
  fields' = \{ fplan \mapsto 0, fmesg \mapsto 0 \} \\
  planSource' = srcManual \\
  lastTool' = 0 \\
  beat' = 0 \\
  curRev' = 0 \\
  now' = 0 \\
  stamp' = 0 \\
  hasStamp' = zfalse \\
  phase' = hbIdle \\
  seenFields' = \{ fplan \mapsto 0, fmesg \mapsto 0 \} \\
  seenSource' = srcManual \\
  seenTool' = 0 \\
  seenRev' = 0 \\
  seenAt' = 0 \\
  seenHas' = zfalse \\
  tries' = 0 \\
  intended' = \{ fplan \mapsto 0, fmesg \mapsto 0 \} \\
  intendedSource' = srcManual \\
  peak' = 0 \\
  manualLive' = zfalse \\
  manualPlan' = 0 \\
  beatHas' = zfalse \\
  beatAt' = 0
\end{schema}

%%----------------------------------------------------------------------
\section{Operations}
%%----------------------------------------------------------------------

\begin{schema}{Tick}
  \Delta Env \\
  \Xi Row \\
  \Xi Heartbeat \\
  \Xi Intent
  \where
  now < maxNow \\
  now' = now + 1 \\
  stamp' = stamp \\
  hasStamp' = hasStamp
\end{schema}

Time passes.

\begin{schema}{ToolCall}
  \Delta Env \\
  \Xi Row \\
  \Xi Heartbeat \\
  \Xi Intent
  \where
  now' = now \\
  stamp' = now \\
  hasStamp' = ztrue
\end{schema}

A tool call ends and the hook touches the stamp file: its modification time is
now.

\begin{schema}{HeartbeatRead}
  \Xi Env \\
  \Xi Row \\
  \Delta Heartbeat \\
  \Xi Intent
  \where
  phase = hbIdle \\
  tries < maxTries \\
  phase' = hbRead \\
  seenFields' = fields \\
  seenSource' = planSource \\
  seenTool' = lastTool \\
  seenRev' = curRev \\
  seenAt' = stamp \\
  seenHas' = hasStamp \\
  tries' = tries
\end{schema}

\texttt{kv.get} of the row, with the stamp read at the start of the tick.
Enabled while the tick has tries left.

\begin{schema}{WriteMesg}
  \Xi Env \\
  \Delta Row \\
  \Xi Heartbeat \\
  \Delta Intent \\
  v? : 0 \upto 1
  \where
  curRev < maxRev \\
  fields' = fields \oplus \{ fmesg \mapsto v? \} \\
  planSource' = planSource \\
  lastTool' = lastTool \\
  beat' = beat \\
  curRev' = curRev + 1 \\
  intended' = intended \oplus \{ fmesg \mapsto v? \} \\
  intendedSource' = intendedSource \\
  peak' = peak \\
  manualLive' = manualLive \\
  manualPlan' = manualPlan \\
  beatHas' = beatHas \\
  beatAt' = beatAt
\end{schema}

A writer other than the heartbeat sets the \texttt{mesg} field.  It runs
whenever it likes, in particular while a heartbeat is in $hbRead$.

\begin{schema}{BiffCall}
  \Xi Env \\
  \Delta Row \\
  \Xi Heartbeat \\
  \Delta Intent
  \where
  curRev < maxRev \\
  fields' = fields \\
  planSource' = planSource \\
  lastTool' = now \\
  beat' = beat \\
  curRev' = curRev + 1 \\
  intended' = intended \\
  intendedSource' = intendedSource \\
  peak' = ( \IF now > peak \THEN now \ELSE peak ) \\
  manualLive' = manualLive \\
  manualPlan' = manualPlan \\
  beatHas' = beatHas \\
  beatAt' = beatAt
\end{schema}

A biff tool handler runs (\texttt{update\_current\_session}): it writes the idle
clock as now, through its own read-modify-write of the row, and takes a
revision.  A heartbeat that read the row before it holds an older stamp than
the call it must not undo.

\begin{schema}{WritePlan}
  \Xi Env \\
  \Delta Row \\
  \Xi Heartbeat \\
  \Delta Intent \\
  v? : 0 \upto 1 \\
  src? : SRC
  \where
  curRev < maxRev \\
  \lnot ( src? = srcAuto \land planSource = srcManual \land fields~fplan \neq 0 ) \\
  fields' = fields \oplus \{ fplan \mapsto v? \} \\
  planSource' = src? \\
  lastTool' = lastTool \\
  beat' = beat \\
  curRev' = curRev + 1 \\
  intended' = intended \oplus \{ fplan \mapsto v? \} \\
  intendedSource' = src? \\
  peak' = peak \\
  manualLive' = ( \IF src? = srcManual \THEN
    ( \IF v? \neq 0 \THEN ztrue \ELSE zfalse ) \ELSE manualLive ) \\
  manualPlan' = ( \IF src? = srcManual \THEN v? \ELSE manualPlan ) \\
  beatHas' = beatHas \\
  beatAt' = beatAt
\end{schema}

A plan write.  The guard is the rule: a plan with $source = auto$ is refused
when the stored $plan\_source$ is manual and the stored plan is not empty.  Any
other write succeeds, replacing the plan and its source, and takes a revision.

\begin{schema}{PlanRefused}
  \Xi PresenceState \\
  v? : 0 \upto 1 \\
  src? : SRC
  \where
  src? = srcAuto \\
  planSource = srcManual \\
  fields~fplan \neq 0
\end{schema}

The refused write: nothing changes, not the plan, not its source, not the
revision.  (The reply says ``unchanged''; the model has no reply.)

\begin{schema}{HeartbeatCommit}
  \Xi Env \\
  \Delta Row \\
  \Delta Heartbeat \\
  \Delta Intent
  \where
  phase = hbRead \\
  curRev = seenRev \\
  curRev < maxRev \\
  fields' = seenFields \\
  planSource' = seenSource \\
  lastTool' = ( \IF seenHas = ztrue \land seenAt > seenTool \THEN seenAt
    \ELSE seenTool ) \\
  beat' = (beat + 1) \mod (maxBeat + 1) \\
  curRev' = curRev + 1 \\
  phase' = hbIdle \\
  seenFields' = seenFields \\
  seenSource' = seenSource \\
  seenTool' = seenTool \\
  seenRev' = seenRev \\
  seenAt' = seenAt \\
  seenHas' = seenHas \\
  tries' = 0 \\
  intended' = intended \\
  intendedSource' = intendedSource \\
  peak' = ( \IF lastTool' > peak \THEN lastTool' \ELSE peak ) \\
  manualLive' = manualLive \\
  manualPlan' = manualPlan \\
  beatHas' = seenHas \\
  beatAt' = seenAt
\end{schema}

\texttt{kv.update(key, row, last=entry.revision)}.  The guard $curRev =
seenRev$ is the compare: the bucket accepts the write only if no write has
landed since the read.  The row written is the row read with two things
moved: the heartbeat counter, and, when a stamp exists, the idle clock to the
later of its value and the stamp.  The plan, its source and every other field
are the ones read, $fields' = seenFields$ and $planSource' = seenSource$; the
guard is what makes that equal to the row now in the bucket.  With no stamp the
idle clock is $seenTool$, unchanged.  The counter wraps at $maxBeat$ only to
keep ProB's state space finite.

\begin{schema}{HeartbeatConflict}
  \Xi Env \\
  \Xi Row \\
  \Delta Heartbeat \\
  \Xi Intent
  \where
  phase = hbRead \\
  curRev \neq seenRev \\
  phase' = hbIdle \\
  seenFields' = seenFields \\
  seenSource' = seenSource \\
  seenTool' = seenTool \\
  seenRev' = seenRev \\
  seenAt' = seenAt \\
  seenHas' = seenHas \\
  tries' = tries + 1 \\
  tries < maxTries
\end{schema}

\texttt{KeyWrongLastSequenceError}: the row changed since the read.  Nothing
is written; the loop goes round again, and $HeartbeatRead$ is no longer
enabled once $tries = maxTries$ (``Session changed concurrently, skip
heartbeat'').

\begin{schema}{NextTick}
  \Xi Env \\
  \Xi Row \\
  \Delta Heartbeat \\
  \Xi Intent
  \where
  phase = hbIdle \\
  tries > 0 \\
  tries' = 0 \\
  phase' = hbIdle \\
  seenFields' = seenFields \\
  seenSource' = seenSource \\
  seenTool' = seenTool \\
  seenRev' = seenRev \\
  seenAt' = seenAt \\
  seenHas' = seenHas
\end{schema}

The next heartbeat tick starts with a fresh budget.  The model lets a tick
end after one conflict, which the code does not (it retries within the tick);
the model has more behaviours than the code, and a safety invariant that holds
in it holds for the code.

\begin{schema}{Quiesce}
  \Xi PresenceState
\end{schema}

The bounds on $curRev$ and $now$ end the history after a few steps.  $Quiesce$
is the stutter that stands for ``nothing more happens''; without it ProB
reports the end of the bounded history as a deadlock, which says something
about the bound and nothing about the protocol.

%%----------------------------------------------------------------------
\section{Model Checking}
%%----------------------------------------------------------------------

\begin{enumerate}
  \item \textbf{The properties}: $PresenceInv$ is a state invariant.  ProB checks
    it over the whole bounded state space.
  \item \textbf{Liveness of the tick}: $EF\; beat \neq 0$: a heartbeat can
    land.  A guard that were never enabled would make the first property
    vacuous.
  \item \textbf{The race is real}: $EF\;(phase = hbRead \land curRev \neq
    seenRev)$.
  \item \textbf{The heartbeat itself moves the idle clock}: $EF\;(phase =
    hbIdle \land beatHas = ztrue \land beatAt > seenTool \land lastTool \geq
    beatAt)$.  $seenTool$ is the idle clock as the heartbeat read it and $beatAt$
    the stamp it committed.  A call of biff before the read is already in
    $seenTool$, and one between the read and the commit makes $curRev \neq
    seenRev$ so the commit does not happen; so the clock reached $beatAt$ by the
    heartbeat alone, not by $BiffCall$.
  \item \textbf{The guard refuses}: $EF$ of an auto plan write that finds a
    non-empty manual plan is the operation $PlanRefused$ enabled, stated as the
    goal $planSource = srcManual \land fields~fplan \neq 0$.
\end{enumerate}

Four frozen variants are the negative controls, each the shipped file with one
thing taken out, each reporting an invariant violation.
\begin{itemize}
  \item \texttt{presence-heartbeat-blindput.tex}: $HeartbeatCommit$ without the
    $curRev = seenRev$ guard (the \texttt{kv.put} the code used to make).  The
    plan written between the read and the write is overwritten: $HeartbeatRead
    \to WritePlan \to HeartbeatCommit$.
  \item \texttt{presence-heartbeat-nostamp.tex}: $HeartbeatCommit$ leaves the idle
    clock alone whatever the stamp (DES-097's bug).  A tool call, time passing,
    and a heartbeat leave $lastTool$ older than the stamp the heartbeat carried:
    property 3.
  \item \texttt{presence-heartbeat-stampback.tex}: $HeartbeatCommit$ sets the idle
    clock to the stamp unconditionally.  A stale stamp read after a later
    biff call moves $lastTool$ backwards: property 2.
  \item \texttt{presence-heartbeat-noguard.tex}: $WritePlan$ without the guard.
    An automatic plan replaces a non-empty manual one: property 4.
\end{itemize}

%%----------------------------------------------------------------------
%% Reproducing the checks.  Run from the repo root; probcli is not assumed to
%% be on PATH.  Bounded run (one probcli at a time):
%%
%%   PROBCLI="${PROBCLI:-$HOME/Applications/ProB/probcli}"
%%   P="-p DEFAULT_SETSIZE 1 -p MAXINT 6 -p MAX_INITIALISATIONS 20 \
%%      -p MAX_OPERATIONS 500 -p TIME_OUT 600000"
%%   ( ulimit -v 8000000; "$PROBCLI" docs/specs/presence-heartbeat.tex $P -model_check )
%%   "$PROBCLI" docs/specs/presence-heartbeat.tex $P -model_check -goal "beat/=0"
%%   "$PROBCLI" docs/specs/presence-heartbeat.tex $P -model_check \
%%       -goal "phase=hbIdle & beatHas=ztrue & beatAt>seenTool & lastTool>=beatAt"
%%
%% Negative controls, each of which must report an invariant violation:
%%
%%   for V in blindput nostamp stampback noguard; do
%%     "$PROBCLI" docs/specs/presence-heartbeat-$V.tex $P -model_check
%%   done
%%----------------------------------------------------------------------

\end{document}
